\subsection{Computational Security}

A scheme is $(t, \epsilon)$-secure if any 
adversary running for time at most $t$ succeeds 
in breaking the scheme with probability at most 
$\epsilon$. However, this definition makes 
no guarantees for after $t$; perhaps the probability
of an attacker succeeding rises sharply after $t$. 

A private-key encryption scheme with key length 
$n$ and for an adversary running
for time $t$, succeeds in breaking the scheme 
with probability at most $c \frac{t}{2^n}$ (for
some constant $c$).

A better approach to security is to say a scheme 
is secure if any \emph{efficient} adversary succeeds in breaking the scheme with at
most \emph{negligible} probability. 

A function $f$ (from natural numbers to non-negative real numbers) 
is polynomially bounded (or polynomial) if there is a constant $c$ 
such that $f(n) < n^c$ for all $n$.

An algorithm $A$ runs in polynomial time if there is a 
polynomial $p$ s.t. for all inputs
$x \in \{0, 1\}^*$, $A(x)$ terminates within $p(|x|)$ steps.

An algorithm is efficient if it runs in polynomial time. 

A function $v(n)$ is negligible if for every $c$, there exists some 
$N$ such that for all $n > N$, $v(n) \leq \frac{1}{n^c}$. 

A scheme is computationally secure if for every 
probabilistic polynomial-time (PPT) adversary $A$
carrying out an attack, the probability that $A$ 
succeeds is negligible. 